\section{Background}
\label{sec:rs-background}

The Recursive Spawning Protocol emerges from four converging research traditions: agent lifecycle management in multi-agent AI systems, accountability and provenance in distributed systems, the structural properties of delegation chains, and hierarchical task decomposition. Each tradition exposes a distinct inadequacy in prior work that RSP is designed to address. This section traces each tradition, establishes the specific gap it reveals, and positions the BX3 Framework as the minimal sufficient architectural substrate for RSP's guarantees.

% --------------------------------------------------------------
\subsection{Agent Lifecycle Management in Multi-Agent Systems}
\label{sec:rs-agent-lifecycle}

Agent lifecycle management encompasses the mechanisms by which agent instances are created, provisioned with authority, execute their delegated scope, transition between states, and are terminated or archived. In single-agent deployments, lifecycle management is largely sequential and bounded by design: a principal issues a directive, the agent executes within a defined scope, and the lifecycle terminates when the task is resolved. Multi-agent systems fundamentally complicate this model because the lifecycle of any given agent may encompass the creation of subordinate agents---recursive delegation---whose own lifecycles must be coherently managed within the same accountability architecture~\cite{intelligent_ai_delegation}.

Recursive delegation has been identified as the dominant operational mode in complex task environments. Park \textit{et al}. formalize the recursion condition: an agent that has accepted a task may further decompose and delegate subtasks, and any subtask may itself be decomposed and delegated further, until a leaf agent reaches a task that is resolvable without further delegation~\cite{arxiv2602.11865}. The resulting structure is a delegation chain---a directed acyclic graph rooted at the original human principal---in which each node carries some portion of the original directive's authority, refined and constrained to its assigned subtask. As chains deepen, the connection between any given node and the human principal at the root becomes attenuated. The node's operating scope is only as well-defined as the delegation chain that provisioned it; if the chain is mutable, the node's scope is mutable; if the chain is opaque, the node's accountability is opaque~\cite{zylos_hierarchical_coordination}.

Agent lifecycle management in production multi-agent systems must address not only the lifecycle of individual agents but the lifecycle of the chain itself. This includes: provisioning new nodes with scoped authority derived from the parent node's refined operating range; tracking the complete ancestry of each node back to a human principal; detecting lifecycle anomalies such as unauthorized re-delegation, scope expansion beyond provisioned bounds, or spawning events that lack a valid parent chain; and terminating nodes cleanly with full forensic preservation. Systems that address lifecycle management only at the node level---without an architectural mechanism for chain-level accountability---leave the critical dependency between parent and child unverified and mutable~\cite{multi_agent_coordination_strategies}.

The enterprise shift toward hierarchical agent teams reflects the operational reality that complex multi-step tasks require layered decomposition in which manager agents delegate to specialized subordinate agents, each operating with reduced context requirements and narrower scope~\cite{ibm_hierarchical_agent_hierarchies}. This pattern makes lifecycle management more critical with every additional layer: each new delegation relationship must be both authorized and auditable, yet without structural enforcement of the delegation chain, hierarchical agent architectures become opaque trees in which the original principal's intent is progressively distorted as it propagates downward.

% --------------------------------------------------------------
\subsection{Accountability and Provenance in Distributed Systems}
\label{sec:rs-accountability}

Provenance in distributed systems is defined as the complete record of the causal chain connecting an output to its input sources and the agents that transformed them~\cite{immutable_audit_trails_blockchain}. Provenance is the structural prerequisite for accountability: without a complete, tamper-evident record of what was done, by whom, under whose authority, and in response to what input, accountability is a stated intention rather than a verifiable property. The distributed systems literature has established that immutable audit trails are the minimal sufficient mechanism for provenance, because mutability of the record is sufficient to destroy the evidentiary chain~\cite{blockchain_ai_governance}.

In AI systems specifically, on-chain provenance for training data, model weights, and inference outputs has been proposed as the primary mechanism for establishing model accountability~\cite{ accountable_ai_dag}. The core argument is that AI outputs are the result of transformations applied to inputs, and the complete causal graph of these transformations must be recorded if the output is to be attributable to its sources. However, the AI provenance literature has largely focused on data provenance within single-agent inference pipelines, not on the delegation provenance across recursive multi-agent chains. These are structurally distinct problems: data provenance traces material flows; delegation provenance traces authority flows. An agent that spawns a child agent is not transforming data---it is authorizing a new actor to act on its behalf. The accountability requirement in delegation provenance is therefore stronger: not only must the chain of transformations be recorded, but the chain of authorizations must also be recorded, and the authorization at each link must be verified as valid before the child node is provisioned.

The transitivity of accountability through delegation chains presents a fundamental verification challenge. If agent $A$ delegates to agent $B$, and $B$ delegates to agent $C$, is $A$ accountable for $C$'s actions? The answer depends on whether the delegation was authorized, whether the scope of $C$'s actions falls within the scope of $A$'s original directive, and whether the chain connecting $A$ to $C$ is intact and unmodified. If the delegation chain is mutable, the transitivity claim cannot be structurally verified; it is a policy assertion, not an architectural guarantee.

% --------------------------------------------------------------
\subsection{Fixed and Mutable Delegation Structures}
\label{sec:rs-delegation}

The structural properties of delegation chains determine what accountability guarantees are achievable. A \textit{mutable} delegation chain is one in which the parent pointer, operating scope, or chain topology of any node can be modified after provisioning---by the node itself, by a parent, by an administrator, or by any actor with sufficient privilege. A \textit{fixed} delegation chain is one in which these properties are established exactly once at provisioning and cannot be modified by any actor thereafter. The distinction is not merely a matter of policy strength; it is a structural classification that determines the categories of failure modes that can arise~\cite{bx3_framework_2026}.

Mutable delegation chains introduce three structurally distinct failure modes that have no architectural remedy within the mutable chain itself. First, \textit{retroactive scope expansion}: a node modifies its own operating scope to encompass actions that were not authorized by its parent, without the parent's knowledge or consent. This is possible whenever the node has the privilege to modify its own configuration. Second, \textit{parent pointer corruption}: the chain topology is altered after the fact, either by a man-in-the-middle attack on the communication channel between nodes, or by a node with administrative privilege modifying the parent pointer of a downstream node. Third, \textit{chain inversion}: the delegation direction is reversed---a subordinate node presents itself as the principal, or a node claims authority over its parent---which is the structural expression of impersonation in the delegation chain~\cite{agent_delegation_vs_impersonation}.

These three failure modes are not addressed by stronger policy, more detailed logging, or better monitoring. They are structural failures that arise from the mutability of the chain itself. A policy that prohibits retroactive scope expansion is effective only if it is architecturally enforced; if the node can modify its own configuration, the policy is unenforceable without structural constraint. Similarly, logging cannot prevent parent pointer corruption; only immutability of the pointer can do that. The conclusion is structural, not empirical: no mutable delegation chain can provide the accountability guarantees required for autonomous AI agents operating in consequential environments, regardless of the policy, monitoring, or cryptographic sophistication applied at the node level.

Fixed delegation chains---in which the parent pointer, scope, and chain topology are established exactly once at provisioning and are immutable thereafter---prevent all three failure modes by construction. The chain is tamper-evident: any modification to the chain topology after provisioning is detectable because the immutable record will not match the modified state. The structural prerequisite for fixed delegation chains is an append-only, cryptographically sealed forensic record maintained independently of the runtime state of any individual node.

% --------------------------------------------------------------
\subsection{Hierarchical Task Decomposition in AI}
\label{sec:rs-task-decomposition}

Hierarchical task decomposition is the process by which a complex, high-level directive is recursively refined into simpler subtasks that can be executed by specialized agents or human workers. The decomposition terminates when all leaf tasks are resolvable without further delegation. This pattern is foundational in classical planning~\cite{hybrid_llm_agents}, in robotics task graphs~\cite{robot_task_decomposition}, and increasingly in LLM-based multi-agent systems where no single model can efficiently execute the full scope of a complex directive~\cite{multi_agent_decomposition_survey}.

Hierarchical decomposition creates delegation chains as a byproduct of task assignment. When a parent node decomposes a task and delegates a subtask to a child node, it establishes a delegation relationship that must be recorded, preserved, and verifiable. In mutable delegation structures, this relationship is established as a policy assertion at delegation time but is not architecturally enforced thereafter. In the BX3 Framework, the Purpose Layer enforces the refinement condition at every spawn event: a new node may only be provisioned if its operating scope is a descendant refinement of the parent node's operating scope, and this condition is verified by the Purpose Layer before the Fact Layer records the immutable parent pointer. The enforcement is structural, not advisory~\cite{bx3_framework_2026}.

The termination condition for hierarchical decomposition is equally important for accountability. Decomposition that never terminates---or that terminates only when the system's resources are exhausted---is a structural failure mode. The RSP enforces termination through a Purpose Layer-defined maximum depth bound: decomposition terminates either when all leaf tasks are resolvable without further delegation or when the escalation threshold is reached, at which point the Bailout Protocol ensures that the unresolved remainder is escalated to the human anchor rather than allowed to propagate indefinitely through the machine-only execution graph~\cite{bailout_protocol_2026}.

Existing approaches to hierarchical task decomposition in multi-agent AI systems do not simultaneously enforce scope refinement, chain immutability, and depthbounded termination. Systems that emphasize decomposition efficiency~\cite{multi_agent_decomposition_survey} typically treat the delegation chain as an implementation detail rather than an accountability mechanism. Systems that address accountability typically do so at the node level, not at the chain level, and therefore cannot prevent the three failure modes identified in Section~\ref{sec:rs-delegation}.

% --------------------------------------------------------------
\subsection{The BX3 Framework as Immutable Parent Chain Substrate}
\label{sec:rs-bx3-substrate}

The BX3 Framework provides the three-layer accountability architecture that serves as the structural substrate for the Recursive Spawning Protocol. The framework defines three orthogonal layers---the Purpose Layer, the Bounds Engine, and the Fact Layer---that operate in mandatory concert: the Purpose Layer authorizes, the Bounds Engine constrains, and the Fact Layer records. No layer can substitute for another, and no layer's function can be subsumed by either of the others without collapsing at least one of the RSP's correctness properties~\cite{bx3_recursive_spawning_integration}.

The \textbf{Purpose Layer} is responsible for mission coherence and drift prevention. At every spawn event, the Purpose Layer validates that the child node's operating scope is a descendant refinement of the parent node's operating scope---never an expansion. This validation is non-negotiable and architectural: the spawn event is not recorded by the Fact Layer until the Purpose Layer has verified the refinement condition. The Purpose Layer also defines the maximum depth bound for the spawn chain, enforcing the termination property by ensuring that no chain can exceed the depth limit without triggering the escalation pathway.

The \textbf{Bounds Engine} defines and enforces the operational envelope of each node. It is the mechanism by which scope constraints are translated into runtime behavior: a node operating within its provisioned bounds is permitted to act; a node attempting to act outside its bounds triggers the Bounds Engine's exception propagation pathway, which leads to either local resolution within the chain or escalation to the human anchor via the Bailout Protocol. The Bounds Engine is not a static permission matrix; it is an active constraint-satisfaction engine that monitors the node's operational state in real time and enforces the envelope as a structural property, not as a policy recommendation.

The \textbf{Fact Layer} maintains the Forensic Ledger---an append-only, cryptographically sealed record of every state transition, every spawn event, every scope refinement, and every exception propagation across the agent network. The Fact Layer is the system of record for chain verification: the immutable parent pointer from every child node to its parent is recorded exactly once at provisioning and cannot be modified thereafter. Chain verification---the process of confirming that the delegation chain from any node to its human anchor is intact, unmodified, and still within its authorized depth bound---proceeds by traversing the parent pointers recorded in the Forensic Ledger, not by querying the nodes themselves. This separation of record from runtime state is architecturally critical: the ledger is authoritative even when the node network is in a degraded or adversarial state.

The BX3 Framework's three-layer architecture is the minimal sufficient structure for RSP's guarantees. No proper subset of the three layers achieves the drift prevention, chain integrity, and termination correctness properties simultaneously. The Purpose Layer alone cannot enforce runtime constraints or maintain the forensic record; the Bounds Engine alone cannot authorize spawn events or enforce refinement constraints; the Fact Layer alone cannot prevent drift or enforce depth bounds. Only the three layers operating in mandatory concert provide the structural guarantee that the delegation chain is immutable, auditable, scope-refined, depth-bounded, and verified to terminate at a human principal~\cite{bx3_framework_2026}.
